% ============================================================
% METHODS
% ============================================================

\section{Methods}
\label{sec:methods}

This section describes our methodology for applying deep reinforcement learning to cryptocurrency portfolio management. We first formalize the problem as a Markov Decision Process (\Cref{sec:problem_formulation}), then detail our environment design including data processing and constraint handling (\Cref{sec:environment}). We present four RL agent architectures (\Cref{sec:agents}) and three baseline strategies (\Cref{sec:baselines}) that operate under identical constraints for fair comparison.

% ------------------------------------------------------------
\subsection{Problem Formulation}
\label{sec:problem_formulation}
% ------------------------------------------------------------

We formulate cryptocurrency portfolio management as a Markov Decision Process (MDP) defined by the tuple $(\mathcal{S}, \mathcal{A}, \mathcal{P}, \mathcal{R}, \gamma)$, where $\mathcal{S}$ is the state space, $\mathcal{A}$ is the action space, $\mathcal{P}$ defines transition dynamics, $\mathcal{R}$ is the reward function, and $\gamma \in [0,1]$ is the discount factor.

\paragraph{Transition Dynamics.} The transition $\mathcal{P}(s_{t+1}|s_t, a_t)$ is determined by two components: (1) exogenous market price movements, which we do not model or predict, and (2) endogenous portfolio mechanics, where the new weights $\vect{w}_t$ become the ``previous weights'' in the next state. We adopt a model-free approach---agents learn directly from environment interactions without estimating $\mathcal{P}$.

\paragraph{State Space.} At each trading day $t$, the state $s_t$ consists of (1) a per-asset historical observation tensor $\mat{X}_t \in \R^{A_t \times 4 \times 60}$ containing normalized OHLCV features over a 60-day lookback window, and (2) the previous portfolio weights $\vect{w}_{t-1} \in \Delta^{A_{t-1}}$, where $A_t$ denotes the number of tradable assets at time $t$. The inclusion of previous weights, analogous to the Portfolio Vector Memory (PVM) mechanism~\citep{jiang2017eIIE}, enables the agent to internalize turnover costs.

\paragraph{Action Space.} The action $a_t$ specifies a target portfolio allocation $\vect{w}_t \in \Delta^{A_t}$ over the current tradable assets. For value-based methods, we discretize this space into 70 delta-based rebalancing actions (\Cref{sec:dqn}); for policy gradient methods, we output continuous weights directly (\Cref{sec:reinforce}).

\paragraph{Constraints.} All portfolios must satisfy:
\begin{align}
    w_t^{(i)} &\geq 0 \quad \forall i \quad \text{(long-only)} \label{eq:long_only} \\
    \sum_{i=1}^{A_t} w_t^{(i)} &= 1 \quad \text{(fully invested, no cash)} \label{eq:simplex} \\
    \|\vect{w}_t - \vect{w}_{t-1}\|_1 &\leq \tau \quad \text{(turnover limit, } \tau = 0.30\text{)} \label{eq:turnover}
\end{align}

The fully-invested constraint~\eqref{eq:simplex} forces the agent to maintain exposure to cryptocurrency risk rather than trivially holding cash, ensuring fair comparison against crypto benchmarks~\citep{jiang2017eIIE, lucarelli2020dqlcrypto}.

\paragraph{Reward Function.} The agent receives a reward based on net portfolio return after transaction costs:
\begin{equation}
    r_t = \log\left(1 + \vect{w}_t^\top \vect{R}_{t+1}\right) - c \cdot \|\vect{w}_t - \vect{w}_{t-1}\|_1
    \label{eq:reward}
\end{equation}
where $\vect{R}_{t+1} \in \R^{A_t}$ is the vector of simple returns from day $t$ to $t+1$, and $c = 0.001$ is the transaction cost coefficient (0.1\% per unit turnover). Critically, the forward returns $\vect{R}_{t+1}$ are never observed by the agent at decision time---they are used only by the environment to compute realized rewards after the action is taken.

\paragraph{Objective.} The agent seeks to maximize expected cumulative discounted rewards:
\begin{equation}
    \max_\pi \; \E_\pi \left[ \sum_{t=0}^{T} \gamma^t r_t \right]
\end{equation}
We use $\gamma = 0.93$ for production training, selected via Bayesian hyperparameter optimization (\Cref{sec:hyperparameters}). This relatively high discount factor indicates that longer-horizon planning is beneficial even with daily transaction costs, though the search space included values as low as 0.5.

% ------------------------------------------------------------
\subsection{Environment Design}
\label{sec:environment}
% ------------------------------------------------------------

Our environment processes over seven years of cryptocurrency market data and enforces realistic trading constraints. We describe the data pipeline, state representation, and constraint enforcement mechanisms.

\subsubsection{Data Pipeline}

We use daily OHLCV (Open, High, Low, Close, Volume) data for liquid cryptocurrencies from 2018-07-01 through 2025-10-31. The tradable universe is defined by a market-cap-weighted index with monthly constituent rebalancing: at the end of month $m-1$, we freeze the index membership for all days in month $m$, preventing intramonth look-ahead bias.

\paragraph{Data Quality.} We apply tiered gap repair rules to handle missing data. Single-day gaps are addressed through forward-filling, propagating the previous day's bar. Gaps of two to five days receive linear interpolation for price channels and log-space interpolation for volume to preserve multiplicative scaling properties. Gaps exceeding five days trigger asset ineligibility---the asset is excluded from the tradable universe until 60 consecutive days of clean data become available, ensuring sufficient history for feature construction.

\paragraph{Cold-Start Eligibility.} An asset is tradable on day $t$ only if it (1) is included in the current month's index membership, and (2) has at least 60 consecutive days of clean OHLCV data immediately prior to $t$. This prevents allocation to newly listed or illiquid assets with insufficient trading history~\citep{jiang2017eIIE, ye2020sarl}.

\subsubsection{State Representation}

For each tradable asset $i$ at time $t$, we construct a feature tensor from the trailing 60-day window:

\paragraph{Price Normalization.} Close, high, and low prices are divided by the asset's closing price on day $t$, so the most recent close equals 1.0 and prior prices are relative. This improves stationarity and has been shown to stabilize training in crypto portfolio agents~\citep{jiang2017eIIE}.

\paragraph{Volume Normalization.} We apply $\log(1 + \text{volume})$, then z-score normalize within each asset's 60-day window, clipping extreme values to $[-5, 5]$.

The resulting observation tensor $\mat{X}_t \in \R^{A_t \times 4 \times 60}$ has dimensions: $A_t$ assets $\times$ 4 channels (close, high, low, volume) $\times$ 60 days. Because $A_t$ varies over time (due to index rebalancing and cold-start rules), we employ canonical padding for value-based methods (\Cref{sec:dqn}).

\paragraph{Universe Changes.} When assets exit the tradable universe (due to index rebalancing or data quality issues), the environment forces liquidation at the last available price and redistributes the weight proportionally across remaining assets, incurring transaction costs. New entrants begin with zero weight and become available for allocation only after satisfying the cold-start requirement.

\subsubsection{Constraint Enforcement}
\label{sec:constraint_enforcement}

The environment supports two constraint handling approaches:

\paragraph{Projection Mode} (used by REINFORCE): If the proposed allocation violates constraints, the environment projects it onto the feasible set using quadratic programming. The agent is trained on the projected allocation without explicit feedback about violations.

\paragraph{Penalty Mode} (used by DQN/DDQN): If the proposed allocation violates constraints, the environment rejects execution and assigns a penalty reward of $-10.0$. The portfolio remains at its previous state, but time advances normally. The penalty magnitude was chosen to be substantially larger than typical single-step rewards (which range from approximately $-0.05$ to $+0.05$), ensuring constraint violations are strongly discouraged. This enables the agent to learn context-dependent feasibility through experience---the same delta action (e.g., ``increase BTC by 10\%'') can be safe or risky depending on current holdings~\citep{lucarelli2020dqlcrypto}.

% ------------------------------------------------------------
\subsection{Agent Architectures}
\label{sec:agents}
% ------------------------------------------------------------

We implement two reinforcement learning agent families: value-based methods (DQN, DDQN) and policy gradient methods (REINFORCE, REINFORCE with learned baseline).

% ----------------------------
\subsubsection{Deep Q-Network (DQN) and Double DQN (DDQN)}
\label{sec:dqn}
% ----------------------------

Our value-based agents address two key challenges: handling variable universe sizes with fixed-dimensional networks, and designing a discrete action space for portfolio rebalancing.

\paragraph{Canonical Padding StateEncoder.}
We assign each of the 37 unique assets a fixed canonical position. For each observation, we populate zero-padded arrays $\mat{X}_{\text{canonical}} \in \R^{37 \times 4 \times 60}$ and $\vect{w}_{\text{canonical}} \in \R^{37}$ at assigned positions, then flatten to 8,917 dimensions and project through a learned linear layer to a 256-dimensional embedding. This preserves asset identity essential for context-dependent Q-values.

\paragraph{Delta-Based Action Catalog.}
We implement 70 discrete portfolio adjustment actions: a hold action; adjust actions that increase/decrease exposure to top-$K$ assets by $\pm5\%$, $\pm10\%$, or $\pm15\%$; rotation actions transferring weight between asset pairs; and rebalancing actions resetting to equal weight or market-cap weight. Each action applies a delta adjustment and renormalizes to the simplex.

\paragraph{Network Architecture and Training.}
The Q-network is a 3-layer MLP (512$\rightarrow$256$\rightarrow$70) with ReLU and dropout. We use experience replay (capacity 10,000), $\epsilon$-greedy exploration decaying from 1.0 to 0.1, and target network updates every 100 steps. DDQN differs only in the TD target: $Q_{\text{target}} = r + \gamma \, Q_{\text{target}}(s', \argmax_{a'} Q_{\text{online}}(s', a'))$, decoupling action selection from evaluation to reduce overestimation bias~\citep{vanHasselt2016ddqn}.

% ----------------------------
\subsubsection{REINFORCE Policy Gradient Agents} 
\label{sec:reinforce} 
% ---------------------------- 
Our policy gradient agents directly parameterize portfolio weights as continuous outputs. \paragraph{Policy Network Architecture.} A GRU layer (input size $4 \times N$, sequence length=60) processes the combined assets' 60-day sequence.  The GRU layer outputs to a fully connected layer where it is concatenated with the previous weights and a binary initialization flag.  This is then fed to another hidden layer then to a final output layer which produces per-asset logits, converted via masked softmax to portfolio weights. We use canonical indexing across 37 positions with a boolean mask for tradable assets. 
 
\paragraph{Training with Dirichlet Sampling.} Actions are sampled from a Dirichlet distribution with concentration parameters $\boldsymbol{\alpha} = \text{softplus}(\text{logits}) / \tau$, where temperature $\tau$ is a parameter introduced to flatten or spike the distribution. When sampled allocations violate constraints, the environment projects them onto the feasible set. Policy gradients use Monte Carlo returns, and the network uses an Adam optimizer.
 
\paragraph{Learned Value Baseline.} To reduce gradient variance, REINFORCE+Baseline adds a value head $V_\phi(s_t)$, which shares the GRU encoder for learning efficiency. The advantage $A_t = G_t - V_\phi(s_t)$ replaces raw returns in the policy gradient. The combined loss is $\mathcal{L} = -\sum_t A_t \log \pi_\theta(a_t | s_t) + 0.5 \sum_t (G_t - V_\phi(s_t))^2$.
 
% ------------------------------------------------------------
\subsection{Baselines}
\label{sec:baselines}
% ------------------------------------------------------------

We implement three baselines operating under identical constraints. \textbf{Equal Weight (EW)} allocates uniformly: $w_t^{(i)} = 1/A_t$. Despite its simplicity, equal weighting has proven competitive by avoiding estimation error~\citep{demiguel2009optimal}. \textbf{Market Cap Weight (MCW)} allocates proportionally to market capitalization, mimicking index funds with minimal turnover. \textbf{Mean-Variance Optimization (MVO)} implements Markowitz optimization~\citep{markowitz1952portfolio}:
\begin{equation}
    \vect{w}_t^* = \argmax_{\vect{w} \in \Delta^{A_t}} \left\{ \vect{w}^\top \hat{\boldsymbol{\mu}}_t - \frac{\lambda}{2} \vect{w}^\top \hat{\boldsymbol{\Sigma}}_t \vect{w} \right\}
\end{equation}
using Ledoit-Wolf shrinkage~\citep{ledoit2004honey} for covariance estimation.
